\section{Partial cylinders and elementary expansions}
\label{sec:partial-cylinders}

Continue to work at a fixed cutoff \(q\).  Fix \(0\le n<q\), and
assume the descending induction hypothesis for every index greater than
\(n\), with arbitrary cellular starting objects.  In this section,
write \(F\) for \(F_q\), and retain \(D_j,\partial D_j,I_j,J_j\)
for their images in \(\mathcal A_q\); an \emph{upper} attachment has index greater than
\(n\).  In particular, Lemma~\ref{lem:upper-cell} and
Propositions~\ref{prop:tree-contractible} and
\ref{prop:whiskered-shape} are available.

Our goal is a functor \(\mathcal P\) with natural endpoint maps
\(p_0,p_1:\mathcal P Y\to Y\).  Each endpoint must lift every
boundary inclusion, while the pair \((p_0,p_1)\) must lift \(J_n\).
Here is how these properties will prove the induction step.  For a
cellular \(X\), the individual lifting property supplies a section
\(h:X\to\mathcal P X\) of \(p_0\), and \(u=p_1h\) is a weak
equivalence.  For an elementary expansion \(i:X\to Y\) with its
unit retraction \(r:Y\to X\), the joint lifting property extends
this comparison over \(Y\) and makes \(iur\) a weak equivalence.
The map \(iu\) is then a retract of \(iur\), so retract closure
and two-out-of-three will give the desired weak equivalence \(i\).
We construct the representing cylinder objects first, then interpret
the operations of \(\mathcal C\) to obtain this functor and its
endpoint properties.

\subsection{The cylinder objects}

Representing path cells by maps out of a coglobular cylinder follows
the strategy of Lanari
\cite[Definitions~5.2--5.3 and~6.1]{Lanari2020}.
Here the cylinders depend on the fixed cutoff and expansion index;
their upper completions will support the descending induction and
the interpretation of every operation of the chosen coherator.

The object \(P_j\) will represent a cylinder on a \(j\)-cell:
a map \(P_j\to Y\) specifies that cylinder in \(Y\), with two
endpoint \(j\)-cells.  The endpoints coincide below dimension \(n\).
In dimension \(n\), a cylinder is a single \((n+1)\)-cell joining
them.  In higher dimensions we will attach cells to make these data
compatible along their common faces.

We construct a coglobular diagram \(P_0,\ldots,P_q\), together with
two maps of coglobular diagrams
\[
             e_\epsilon:D_\bullet\longrightarrow P_\bullet,
             \qquad \epsilon=0,1.
\]
For a globular set \(S\) of dimension at most \(q\), put
\begin{equation}\label{eq:cylinder-realization}
 P(S)=\colim_{(d_j\longrightarrow S)}P_j.
\end{equation}
The indexing category is the category of cells and their faces in
\(S\).  In particular, the structure maps into \(P_j\) specify a
map \(P(\partial d_j)\to P_j\), and the endpoint maps induce
\(e_\epsilon(S):F(S)\to P(S)\).  This notation will also be used
while only the cylinder objects needed by \(S\) have been constructed.

Given the objects in dimensions below \(j\), a boundary cylinder and
one prescribed endpoint are represented by \(A_j^\epsilon\); prescribing
both endpoints instead gives \(A_j^{01}\).  Their precise definitions are
\begin{align}
 A_j^\epsilon
   &=P(\partial d_j)\amalg_{\partial D_j,e_\epsilon}D_j,
       \label{eq:cylinder-one-sided}\\
 A_j^{01}
   &=P(\partial d_j)
       \amalg_{\partial D_j\amalg\partial D_j}
                    (D_j\amalg D_j).
       \label{eq:cylinder-two-sided}
\end{align}
In Equation~\eqref{eq:cylinder-two-sided}, the first attaching map uses
both endpoints on the boundary, and the second uses the two boundary
inclusions.  Empty boundaries have their usual meaning at \(j=0\).

The two domains have different tasks.  The two-sided domain records
the boundary and both endpoint values that an operation on cylinders
must preserve.  We will choose \(P_j\) receiving these data.  The
resulting one-sided maps \(A_j^\epsilon\to P_j\) will admit
retractions, so boundary-cylinder data with one prescribed endpoint
extend to a cylinder in every target model.  The joint endpoint lifting problem will use a
separate tree of prescribed data.

For \(j<n\), set \(P_j=D_j\), with the usual coglobular maps and
with both endpoints equal to the identity.  Set
\[
                  P_n=D_{n+1},
\]
with endpoints the source and target \(n\)-faces.  Its boundary map
is the inclusion of their common boundary: all dimensions below \(n\)
are stationary.  Thus
\begin{equation}\label{eq:cylinder-critical-level}
 P(\partial d_n)=\partial D_n,\qquad
 A_n^\epsilon=D_n,\qquad
 A_n^{01}=\partial D_{n+1},
\end{equation}
and \(A_n^{01}\to P_n\) is \(I_{n+1}\).

For \(n<j\le q\), use the ordinary small-object construction on
\(\{I_k:k>n\}\) to factor the terminal map as
\begin{equation}\label{eq:cylinder-upper-completion}
 A_j^{01}\longrightarrow P_j\longrightarrow 1,
\end{equation}
where the first map is upper cellular and the second lifts every
upper boundary inclusion.  This completion may attach infinitely
many cells; the objects \(P_j\) need not be finitely cellular.
The given map from \(P(\partial d_j)\)
and the two endpoint disks supply the coglobular and endpoint maps at
this level.  Their identities hold because they already agree on the
boundary.  These choices complete the construction of \(P_\bullet\).

We record the cellularity needed in its analysis.  Inductively,
\(P(\partial d_j)\) is cellular: apply the globular skeletal
filtration and use the cellular maps
\(P(\partial d_k)\to P_k\) at earlier levels.  For \(j<n\),
the latter map is \(I_j\).  At \(j=n\), it attaches two
\(n\)-cells to \(\partial D_n\), obtaining
\(\partial D_{n+1}\), and then attaches the joining
\((n+1)\)-cell.  For \(j>n\), it attaches the two endpoint
\(j\)-cells, followed by the completion in
Equation~\eqref{eq:cylinder-upper-completion}.  This proves cellularity
of all the boundary objects, the two domains in
Equations~\eqref{eq:cylinder-one-sided}--\eqref{eq:cylinder-two-sided},
and the \(P_j\).  Moreover, for \(j>n\),
\begin{equation}\label{eq:cylinder-upper-one-sided}
                     A_j^\epsilon\longrightarrow P_j
\end{equation}
is upper cellular: it first attaches the missing endpoint \(j\)-cell
and then makes the chosen upper attachments.  The two-sided maps
\(A_j^{01}\to P_j\) are cellular in the range \(n\le j\le q\)
used below.

\subsection{The exact skeletal gluing}

We next deduce contractibility of the prescribed-data objects from the
whiskered shapes of the preceding section.  The following formula
specifies which cylinder objects are used at each stage and keeps track
of all shared faces.

\begin{lemma}\label{lem:cylinder-skeletal-gluing}
Fix \(\epsilon\in\{0,1\}\).  For a finite globular set \(S\),
define, whenever the required cylinder objects have been constructed,
\begin{equation}\label{eq:cylinder-skeletal-stage}
 C_\ell(S)=F(S)\amalg_{F(\sk_\ell S),e_\epsilon}
                                 P(\sk_\ell S).
\end{equation}
For \(\ell>n\), there is a pushout square
\begin{equation}\label{eq:cylinder-skeletal-pushout}
\begin{tikzpicture}[baseline=(current bounding box.center),
                   node distance=1.35cm and 2.5cm,>=Stealth]
\node (a) {\(\displaystyle\coprod_{s\in S_\ell} A_\ell^\epsilon\)};
\node (b) [right=of a] {\(C_{\ell-1}(S)\)};
\node (c) [below=of a] {\(\displaystyle\coprod_{s\in S_\ell} P_\ell\)};
\node (d) at (b |- c) {\(C_\ell(S)\).};
\draw[->] (a) -- (b);
\draw[->] (a) -- (c);
\draw[->] (b) -- (d);
\draw[->] (c) -- (d);
\end{tikzpicture}
\end{equation}
Here \(S_\ell\) is the set of actual \(\ell\)-cells of the
globular set \(S\).  If \(S\) is a pasting scheme, then
\begin{equation}\label{eq:cylinder-skeletal-ends}
 C_n(S)=H_\epsilon(S),\qquad
 C_{j-1}(d_j)=A_j^\epsilon\quad(j>n),\qquad
 C_{\dim(S)}(S)=P(S).
\end{equation}
The first identity uses the whisker orientation for which the retained
copy of \(S\) is endpoint \(\epsilon\).
\end{lemma}

\begin{proof}
The ordinary skeleton \(\sk_\ell S\) is obtained from
\(\sk_{\ell-1}S\) by attaching one globe along its boundary for
each \(s\in S_\ell\).  Both \(F\) and the construction
\(P(-)\) in Equation~\eqref{eq:cylinder-realization} preserve these
colimits.  More explicitly, a map \(C_\ell(S)\to T\) consists of
a map \(F(S)\to T\) and maps from its cylinder objects through
level \(\ell\), compatible on faces and on endpoint
\(\epsilon\).  After specifying the same data through level
\(\ell-1\), extending them at a particular cell \(s\) amounts
to choosing a map \(P_\ell\to T\) with its values already
prescribed on \(P(\partial d_\ell)\) and on the retained
endpoint \(D_\ell\).  Those prescribed values form precisely a
map \(A_\ell^\epsilon\to T\).  This is the universal property
of Equation~\eqref{eq:cylinder-skeletal-pushout}.

In particular, faces shared by two cells use the same map from the
previous skeleton.  No additional copy of a shared face is introduced.
The argument concerns colimits of the specified cylinder objects, so it
also applies when an earlier upper completion has attached cells of
dimension less than its cylinder index.

Through dimension \(n-1\), all the cylinders are stationary.  At
dimension \(n\), each cell acquires one duplicate and one joining
cell; the original higher cells retain their original boundaries.
This proves \(C_n(S)=H_\epsilon(S)\), including the case
\(\dim(S)<n\), in which there are no whiskers.  The other two
identities follow from \(\sk_{j-1}d_j=\partial d_j\) and
\(\sk_{\dim(S)}S=S\), respectively.
\end{proof}

\begin{proposition}\label{prop:cylinder-contractibility}
Every \(P_j\) is contractible.  For \(n<j\le q\), both
one-sided domains \(A_j^\epsilon\) are contractible, and
\(A_j^\epsilon\to P_j\) is an upper cellular weak equivalence.
For every pasting scheme \(S\) of height at most \(q\), the model
\(P(S)\) is cellular and contractible.
\end{proposition}

\begin{proof}
The objects \(P_j\) for \(j\le n\) are free globes and are
contractible by Proposition~\ref{prop:shape-contractible}.  Proceed
by induction on \(j>n\), handling both endpoints.  By
Lemma~\ref{lem:cylinder-skeletal-gluing}, \(A_j^\epsilon\) is
obtained from \(H_\epsilon(d_j)\) by pushing out the maps
\(A_\ell^\epsilon\to P_\ell\) on the boundary cells of
dimensions \(n<\ell<j\), in increasing order of \(\ell\).
At each dimension, attach its finitely many cells one at a time.
The resulting iterated pushout is the coproduct pushout of
Lemma~\ref{lem:cylinder-skeletal-gluing}, by its universal property.
The starting object is cellular and contractible by
Proposition~\ref{prop:whiskered-shape}.  Each map used in this
filtration is already an upper cellular weak equivalence by the
induction on \(j\).  Lemma~\ref{lem:upper-cell} therefore makes
every one of its pushouts a weak equivalence.  All intermediate objects
are cellular, so the lemma applies at every stage.  It follows that
\(A_j^\epsilon\) is contractible.  When \(j=n+1\), this
filtration has no intervening stages:
\(A_{n+1}^\epsilon=H_\epsilon(d_{n+1})\).

The upper cellular map \(A_j^\epsilon\to P_j\) preserves all
cells through dimension \(n\), by Lemma~\ref{lem:high-cells}.
Thus \(P_j\) is nonempty and inherits fillers of dimensions
\(1,\ldots,n\) from \(A_j^\epsilon\).  Its map to the
terminal object lifts the upper boundary inclusions by construction,
so it has every remaining boundary filler.  Hence \(P_j\) is
contractible.  At \(n=0\), the inherited condition is simply
nonemptiness.  Both terminal maps are now weak equivalences, so
\(A_j^\epsilon\to P_j\) is a weak equivalence.  Its upper
cellularity was established in Equation~\eqref{eq:cylinder-upper-one-sided}.
This completes the induction on \(j\).

For a pasting scheme \(S\) with \(\dim(S)\ge n\), start with
\(C_n(S)=H_\epsilon(S)\) and use the same filtration through
all dimensions \(n<\ell\le\dim(S)\).  It ends at \(P(S)\).
Every step is again a cellular weak equivalence by
Lemma~\ref{lem:upper-cell}; its cellular, contractible starting object
therefore gives the assertion.  If \(\dim(S)<n\), then
\(P(S)=F(S)\), so Proposition~\ref{prop:shape-contractible}
applies directly.
\end{proof}

\subsection{Interpreting every operation of the coherator}

Contractibility of \(P(S)\) is the output needed from the geometric
construction.  It will let us extend each operation from its already
prescribed boundary and endpoints.  Making that extension once on the
representing objects will ensure naturality in the model.

For \(Y\in\mathcal A_q\), form the globular set with cells
\begin{equation}\label{eq:partial-path-cells}
           (\mathcal P Y)_j=\Hom_{\mathcal A_q}(P_j,Y),
                     \qquad 0\le j\le q,
\end{equation}
and complete it above dimension \(q\) by unique fillers of all
boundaries.  The coglobular identities give its globular identities.
Precomposition with \(e_\epsilon\) gives globular maps
\(p_\epsilon:\mathcal P Y\to Y\) through dimension \(q\),
which extend uniquely above it.  Equation~\eqref{eq:cylinder-realization}
gives the natural representation
\begin{equation}\label{eq:partial-path-tuples}
 \{\text{globular maps }S\to\mathcal P Y\}
       \cong\Hom_{\mathcal A_q}(P(S),Y)
             \qquad(\dim(S)\le q).
\end{equation}

\begin{proposition}\label{prop:partial-path}
The construction in Equation~\eqref{eq:partial-path-cells} admits
interpretations of every operation of \(\mathcal C\), making
\(\mathcal P:\mathcal A_q\to\mathcal A_q\) a functor and
both \(p_\epsilon\) natural model morphisms.  They lift all
boundary inclusions \(I_j\).  Their joint map
\[
              (p_0,p_1):\mathcal P Y\longrightarrow Y\times Y
\]
lifts \(J_n\).
\end{proposition}

\begin{proof}[Construction of the algebraic interpretation]
Interpret the generators in the \emph{original free construction order}
of \(\mathcal C\).  At each stage, the globular set
\(\mathcal P Y\) already interprets the entire preceding theory,
naturally in the full model \(Y\), and both endpoints preserve that
preceding theory.  This holds initially for the globular structure.

Consider a new generating filler \(\theta\), of output dimension
\(j>0\), on a pasting-scheme arity \(S\).  If \(j>q\),
evaluate its already interpreted boundary and take its unique filler
in \(\mathcal P Y\).  The endpoint equations follow from the
boundary equations and unique fillers in \(Y\).  This is natural
in \(Y\).

Suppose \(j\le q\).  The generator condition gives
\(\dim(S)\le j\), so Equation~\eqref{eq:partial-path-tuples}
applies.  If \(j<n\), the entire tuple lies in stationary
dimensions: \(P(S)=F(S)\) and \(P_j=D_j\).  Use the original
operation of \(Y\).  Its prescribed boundary follows from the
preceding endpoint compatibility, since the endpoint maps are identical
isomorphisms on all the cell sets involved.

It remains to handle \(n\le j\le q\).  View the given operation,
in the fixed category of full models, as a map
\[
                         \theta:D_j\longrightarrow F(S).
\]
This map exists independently of the interpretation being constructed
on \(\mathcal P Y\).  Its required endpoint values are the two
existing maps
\begin{equation}\label{eq:partial-path-endpoint-prescription}
             e_\epsilon(S)\theta:D_j\longrightarrow P(S).
\end{equation}
The two face operations of \(\theta\) belong to the preceding
theory.  Evaluating either one on an \(S\)-tuple gives, naturally
in \(Y\), a function
\[
 \Hom_{\mathcal A_q}(P(S),Y)
       \longrightarrow\Hom_{\mathcal A_q}(P_{j-1},Y).
\]
By Yoneda, these functions are represented by maps
\(P_{j-1}\to P(S)\).  Their parallelism is an equation in the
preceding theory, so they agree on all their common faces and give
\[
                 b_\theta:P(\partial d_j)\longrightarrow P(S).
\]
The previous endpoint compatibility says exactly that
\begin{equation}\label{eq:partial-path-boundary-compatibility}
 b_\theta e_\epsilon(\partial d_j)
       =e_\epsilon(S)(\theta|_{\partial D_j}).
\end{equation}
Consequently \(b_\theta\) and the two maps in
Equation~\eqref{eq:partial-path-endpoint-prescription} glue to a map
\begin{equation}\label{eq:partial-path-two-sided-prescription}
                         A_j^{01}\longrightarrow P(S).
\end{equation}

The map \(A_j^{01}\to P_j\) is cellular for
\(n\le j\le q\), and \(P(S)\) is contractible by
Proposition~\ref{prop:cylinder-contractibility}.  Its terminal map
therefore lifts every boundary inclusion, and hence every cellular map.
It extends Equation~\eqref{eq:partial-path-two-sided-prescription} to
a map \(P_j\to P(S)\).  Precomposition with this chosen map is
the required operation on \(\mathcal P Y\).  Its boundary and
endpoint equations hold by construction.

In this step a face operation can be a composite of output dimension
\(j-1\) whose arity \(S\) has height \(j\).  Its evaluation
and Yoneda representation use the whole preceding theory.  Preservation
of its generators makes the endpoints preserve all their composites
and substitutions as well.  The inequality \(\dim(S)\le j\)
was used for the new generator; no smaller arity bound on its composite
faces is needed.

Make the representing choices once, independently of \(Y\).
Postcomposition by a map \(Y\to Z\) then preserves every
chosen operation.  At limit stages of the free theory, the compatible
interpretations give the interpretation of its union.  Freeness now
provides a model of the full \(\mathcal C\), natural in \(Y\),
and both endpoints are full model morphisms.  The construction leaves
the chosen globular set unchanged; it is therefore still
\(q\)-coskeletal.
\end{proof}

\begin{proof}[Lifting for an individual endpoint]
For \(j\le q\), an \(I_j\)-lifting problem against
\(p_\epsilon\) specifies a boundary cylinder and one endpoint
disk with matching endpoint boundary.  By
Equation~\eqref{eq:cylinder-one-sided}, it is a map
\(A_j^\epsilon\to Y\), and a solution is an extension to
\(P_j\to Y\).

For \(j<n\), the map \(A_j^\epsilon\to P_j\) is the
identity.  For \(j=n\), it is one face inclusion
\(D_n\to D_{n+1}\), which has a retraction given by a unit
operation.  For \(j>n\), it is an upper cellular weak equivalence
by Proposition~\ref{prop:cylinder-contractibility}, and has a
retraction by Lemma~\ref{lem:upper-cell}.  Precomposing with these
retractions gives the desired extensions into every \(Y\).
For \(j>q\), the map \(I_j\) is an isomorphism in
\(\mathcal A_q\).  Thus both endpoints lift all \(I_j\),
and belong to \(\mathcal W\) by Lemma~\ref{lem:lifting-test}.
\end{proof}

\begin{proof}[Lifting for the joint endpoint]
A lifting problem against \(J_n\) consists of a map
\(h:P_n\to Y\) and two maps \(a_0,a_1:D_{n+1}\to Y\)
whose source faces are the two endpoints of \(h\).  Its representing
object is
\begin{equation}\label{eq:partial-path-tree-domain}
 B=P_n\amalg_{D_n\amalg D_n}(D_{n+1}\amalg D_{n+1}),
\end{equation}
where the attaching maps use the two endpoints of \(P_n\) and
the source of each new disk.  Since \(P_n=D_{n+1}\), its
presentation in the depth-\(n\) slot has four \(n\)-cells and
three \((n+1)\)-cells:
\[
 h:x_0\longrightarrow x_1,\qquad
 a_0:x_0\longrightarrow y_0,\qquad
 a_1:x_1\longrightarrow y_1.
\]
For \(n>0\), the four \(n\)-cells have their common lower
globular boundary.  For \(n=0\), they are four ordinary vertices.
The underlying graph is the tree \(y_0-x_0-x_1-y_1\).
Thus \(B\) is cellular and contractible by
Proposition~\ref{prop:tree-contractible}.

\begin{fixedfigure}
\centering
\begin{tikzpicture}[>=Stealth,
    vertex/.style={circle,fill=black,inner sep=1.8pt},
    every node/.style={font=\small}]
\node[vertex,label=above left:\(x_0\)] (x0) at (0,1.7) {};
\node[vertex,label=above right:\(x_1\)] (x1) at (4.2,1.7) {};
\node[vertex,label=below left:\(y_0\)] (y0) at (0,0) {};
\node[vertex,label=below right:\(y_1\)] (y1) at (4.2,0) {};
\draw[->,thick] (x0) -- node[above] {\(h\)} (x1);
\draw[->,thick] (x0) -- node[left] {\(a_0\)} (y0);
\draw[->,thick] (x1) -- node[right] {\(a_1\)} (y1);
\draw[->,thick,dashed] (y0) -- node[above] {\(k\)} (y1);
\end{tikzpicture}
\caption{The solid arrows form the tree representing the joint-endpoint
lifting data \(B\).  Vertices denote \(n\)-cells and arrows denote
\((n+1)\)-cells; all vertices have the same lower boundary when \(n>0\).
Adjoining the dashed comparison \(k\) is one upper boundary attachment
\(I_{n+1}\), giving \(A_{n+1}^{01}\).  The square depicts boundary
data; its higher filling cells belong to the remaining upper completion,
which gives \(P_{n+1}\).}
\label{fig:joint-endpoints}
\end{fixedfigure}


The cylinder on the source \(n\)-face is already \(P_n\).
Attach the missing comparison \(y_0\to y_1\), one
\((n+1)\)-cell, to supply the cylinder on the target face.  Together
with the two endpoint disks this is exactly
\(A_{n+1}^{01}\).  Hence the natural map to \(P_{n+1}\)
factors as
\[
                 B\longrightarrow A_{n+1}^{01}
                        \longrightarrow P_{n+1},
\]
where the first map is one \(I_{n+1}\)-attachment and the second
is the chosen upper completion.  It is an upper cellular map between
cellular contractible objects, so it is a weak equivalence and splits
by Lemma~\ref{lem:upper-cell}.  A retraction
\(P_{n+1}\to B\), followed by the prescribed map \(B\to Y\),
solves the lifting problem.  This also applies when \(q=n+1\):
the missing comparison is a legitimate \(I_q\)-attachment, and
all subsequent required fillers are provided by the same construction.
This completes the proof of Proposition~\ref{prop:partial-path}.
\end{proof}

\subsection{An arbitrary attaching cell}

\begin{proof}[Completion of the descending induction]
Let \(X\) be cellular in \(\mathcal A_q\), let
\(a:D_n\to X\) be arbitrary, and form the algebraic pushout
\[
              Y=X\amalg_{D_n}D_{n+1},\qquad i:X\longrightarrow Y.
\]
The unit on \(a\) gives the elementary retraction
\(r:Y\to X\), with \(ri=\id_X\).  Because
\(p_0:\mathcal P X\to X\) lifts every boundary inclusion and
\(X\) is cellular, there is a map \(h:X\to\mathcal P X\)
with \(p_0h=\id_X\).  Set
\[
                  u=p_1h:X\longrightarrow X,
             \qquad \bar u=ur:Y\longrightarrow X.
\]
Both endpoints are weak equivalences; thus two-out-of-three first gives
\(h\in\mathcal W\) and then \(u\in\mathcal W\).
Also \(\bar u i=u\).

Consider the square
\begin{equation}\label{eq:partial-path-final-square}
\begin{tikzpicture}[baseline=(current bounding box.center),
                   node distance=1.45cm and 3.2cm,>=Stealth]
\node (a) {\(X\)};
\node (b) [right=of a] {\(\mathcal P Y\)};
\node (c) [below=of a] {\(Y\)};
\node (d) at (b |- c) {\(Y\times Y\).};
\draw[->] (a) -- node[above] {\(\mathcal P(i)h\)} (b);
\draw[->] (a) -- node[left] {\(i\)} (c);
\draw[->] (b) -- node[right] {\((p_0,p_1)\)} (d);
\draw[->] (c) -- node[below] {\((\id_Y,i\bar u)\)} (d);
\end{tikzpicture}
\end{equation}
Naturality and \(\bar u i=u\) give its commutativity explicitly:
\[
 (p_0,p_1)\mathcal P(i)h=(i,iu)
                         =(\id_Y,i\bar u)i.
\]
The right-hand map lifts \(J_n\) by
Proposition~\ref{prop:partial-path}, so it lifts its pushout \(i\).
Indeed, restrict the top map along the given \(a:D_n\to X\),
and restrict the bottom map along the new disk
\(D_{n+1}\to Y\).  Solve this \(J_n\)-square; its solution
agrees with the original top map on \(D_n\), and the actual pushout
universal property glues the two maps.  This produces
\(f:Y\to\mathcal P Y\) with
\[
                  p_0f=\id_Y,\qquad p_1f=i\bar u.
\]
This use of the pushout requires no condition on the selected cell
\(a\) beyond being a model morphism from \(D_n\).
Two-out-of-three again gives
\[
                         v:=i\bar u\in\mathcal W.
\]

Finally, \(iu:X\to Y\) is a retract of \(v=iur:Y\to Y\)
in the arrow category.  The maps of arrows in the two directions have
domain maps \(i\) and \(r\), respectively, and codomain map
\(\id_Y\) in both cases.  The identities
\[
                vi=iu,\qquad (iu)r=v,\qquad ri=\id_X
\]
verify the two commuting squares and their identity composite.
Retract closure gives \(iu\in\mathcal W\), and two-out-of-three
with \(u\in\mathcal W\) gives \(i\in\mathcal W\).
Every map in this argument is a morphism of the full algebraic models;
the conclusion therefore includes all freely generated composites and
coherence cells of the pushout, for the arbitrary attaching cell \(a\).

This proves the induction step at \(n\).  Together with the top
case established in the preceding section, descending induction proves
the elementary expansion assertion for every index at this fixed
cutoff \(q\).
\end{proof}

\begin{proof}[Proof of Theorem~\ref{thm:main}]
Return to the original category \(\mathcal A\), and write
\(i:X\to Y\) for the pushout in the theorem.  Fix an instance of
the exact-boundary test in dimension \(k\).  Choose
\(q\ge\max\{n+1,k+1\}\).  By Lemma~\ref{lem:cutoff}, the
reflector \(R_q\) preserves the cellular presentation of \(X\)
and the actual pushout defining \(Y\).  The fixed-cutoff theorem
therefore applies to \(R_q(i)\).

The reflector preserves all cells through dimension \(q\), with
their face maps.  Consequently the given boundary and target disk are
unchanged by reflection.  The source disk and its \((k+1)\)-cell
homotopy supplied by the fixed-cutoff theorem are also in unchanged
dimensions.  They are therefore a source disk with the original exact
boundary and an actual homotopy in \(Y\) before reflection.
Proposition~\ref{prop:weak-test} proves that \(i\) is a weak
equivalence.  This holds for every \(k\), every \(n\ge0\),
every cellular \(X\), and the arbitrary fixed coherator
\(\mathcal C\), as required.
\end{proof}
